
% =============================================================
% APPENDIX
% =============================================================
\appendix

\section{Proofs}
\label{app:proofs}


\subsection{Clopper--Pearson UCB Validity}
\label{app:proof_ucb}

We show that $\hat{R}^{\mathrm{upper}}_{1-\delta}(t)$ in Eq.~\eqref{eq:ucb} satisfies $\Pr(R(t) \leq \hat{R}^{\mathrm{upper}}_{1-\delta}(t)) \geq 1 - \delta$.

\begin{proof}
Under the Bernoulli structure, the observed failure count among $\hat{m}_{\mathrm{cal}}(t)$ selected instances follows $\hat{w}_{\mathrm{cal}}(t) \sim \mathrm{Binomial}(\hat{m}_{\mathrm{cal}}(t),\, R(t))$. Define the CDF of the empirical failure rate:
\begin{equation}
  D(r \mid R(t)) = \Pr\!\left(\hat{R}(t) \leq r \mid R(t)\right),
\end{equation}
where $\hat{R}(t) = \hat{w}_{\mathrm{cal}}(t) / \hat{m}_{\mathrm{cal}}(t)$. By definition:
\begin{equation}
  \hat{R}^{\mathrm{upper}}_{1-\delta}(t) = \sup\!\left\{R \in [0,1] : D(\hat{r}_{\mathrm{cal}}(t) \mid R) \geq \delta\right\}.
\end{equation}
Since $D(r \mid R)$ is monotonically decreasing in $R$ (higher true rate makes low empirical rates less likely), $R(t) > \hat{R}^{\mathrm{upper}}_{1-\delta}(t)$ implies $D(\hat{r}_{\mathrm{cal}}(t) \mid R(t)) < \delta$. Therefore:
\begin{equation}
  \Pr\!\left(R(t) > \hat{R}^{\mathrm{upper}}_{1-\delta}(t)\right)
  \leq \Pr\!\left(D(\hat{r}_{\mathrm{cal}}(t) \mid R(t)) < \delta\right)
  \leq \delta,
\end{equation}
where the last step uses the super-uniformity property: $\Pr(D(\hat{r}_{\mathrm{cal}}(t) \mid R(t)) < \delta) \leq \delta$ for discrete distributions. Hence $\Pr(R(t) \leq \hat{R}^{\mathrm{upper}}_{1-\delta}(t)) \geq 1 - \delta$.
\end{proof}

\subsection{Bernoulli Structure under Two-Mode Routing}
\label{app:proof_routing}

\begin{proof}[Proof of Lemma~\ref{lem:routing}]
For fixed thresholds $(t_1, t_2)$, we define the selection indicator:
\begin{equation}
  I_i := \ind\!\left\{U_1^{(i)} \!\leq\! t_1\right\} + \ind\!\left\{U_1^{(i)} \!>\! t_1 \;\text{and}\; U_2^{(i)} \!\leq\! t_2\right\}.
\end{equation}
Both $U_1^{(i)} = U(\mathcal{J}(\mathbf{x}_i, \hat{y}_i))$ and $U_2^{(i)} = U(\mathcal{J}(\mathbf{x}_i, \hat{y}_i, \mathcal{E}_i))$ are deterministic functions of $(\mathbf{x}_i, \hat{y}_i)$: the judge's computation and the retrieval process $\mathcal{E}_i$ are both determined by the input. Hence $I_i$ is a deterministic function of $(\mathbf{x}_i, \hat{y}_i)$.

The correctness of the selected verdict is:
\begin{align}
  c_i &= \ind\!\left\{U_1^{(i)} \!\leq\! t_1\right\} c_1^{(i)} \nonumber \\
  &\;+ \ind\!\left\{U_1^{(i)} \!>\! t_1 \,\wedge\, U_2^{(i)} \!\leq\! t_2\right\} c_2^{(i)},
\end{align}
which is a deterministic function of $(\mathbf{x}_i, \hat{y}_i, y^*_i)$, since the routing decision depends on $(\mathbf{x}_i, \hat{y}_i)$ and correctness additionally depends on $y^*_i$.

Since $(\mathbf{x}_i, \hat{y}_i, y^*_i)$ are i.i.d.\ from $\mathcal{D}$:
\begin{enumerate}
  \item \textbf{Independence}: The selection events $\{I_i = 1\}$ depend on disjoint instances, so they are mutually independent. Conditioning on selection preserves independence across instances.
  \item \textbf{Identical distribution}: Each selected instance satisfies $(\mathbf{x}_i, \hat{y}_i, y^*_i) \sim \mathcal{D}_{t_1,t_2}$, where $\mathcal{D}_{t_1,t_2} := \mathcal{D} \mid I = 1$ is the conditional distribution.
\end{enumerate}
Note that $I_i \in \{0,1\}$ because the two indicator terms are mutually exclusive: the second requires $U_1^{(i)} > t_1$, which precludes the first.

The failure indicators $W_i = 1 - c_i$ over the selected subset are therefore i.i.d.\ $\mathrm{Bernoulli}(R(t_1, t_2))$, where $R(t_1, t_2) = \mathbb{E}[W \mid I = 1]$.
\end{proof}

\paragraph{Scope and limitations.}
The guarantee is tied to the retrieval snapshot used during calibration. Because retrieval relies on a non-stationary web index, a query may return different evidence at deployment than it did at calibration. Such retrieval drift affects the guarantee only if it changes the conditional error rate of Mode~2 verdicts. As long as this error rate remains stable, the calibrated thresholds continue to transfer. We evaluate this scenario in Section~\ref{sec:drift} by re-querying the web several months after calibration and find that FDR control is preserved. If the quality of the retrieved evidence changes substantially, recalibration restores the guarantee. Likewise, any change to the task, candidate model, or judge model violates the exchangeability assumption and requires recalibration.

\section{Experimental Details}
\label{app:expdetails}

\subsection{Notation}
\label{app:notation}

Table~\ref{tab:notation_full} provides a complete reference of all symbols.

\begin{table}[!ht]
\small
\centering
\caption{Notation reference.}
\label{tab:notation_full}
\begin{tabular}{@{}p{1.7cm}p{5.2cm}@{}}
\toprule
\textbf{Symbol} & \textbf{Description} \\
\midrule
$\mathbf{x},\; \hat{y},\; y^*$ & Question, candidate answer, ground truth \\
$F,\; \mathcal{J}$ & Examinee LLM, judge LLM \\
$v \in \{0,1\}$ & Evaluation verdict \\
$e^* \in \{0,1\}$ & Ground-truth evaluation label \\
$\mathcal{A}(v, e^*)$ & Admission function: $\ind[v\!=\!e^*]$ \\
$\mathcal{A}_{\mathrm{ref}},\; \lambda_A$ & Relevance function and threshold \\
$W$ & Failure indicator: $\ind[v \neq e^*]$ \\
$U,\; U_1,\; U_2$ & Uncertainty (general / Mode 1 / Mode 2) \\
$v_1,\; v_2$ & Verdict from Mode 1 / Mode 2 \\
$c_1,\; c_2$ & Correctness of Mode 1 / Mode 2 \\
$\mathcal{E}$ & Retrieved external evidence \\
$\mathcal{D}_{\mathrm{cal}}$ & Calibration set ($N$ instances) \\
$\hat{m}_{\mathrm{cal}}$ & Selection size on calibration set \\
$\hat{w}_{\mathrm{cal}}$ & Error count on calibration set \\
$\hat{r}_{\mathrm{cal}}$ & Empirical FDR \\
$R(t)$ & True FDR \\
$\hat{R}^{\mathrm{upper}}_{1-\delta}$ & Clopper--Pearson UCB \\
$\hat{t}_1,\; \hat{t}_2$ & Calibrated thresholds (Mode 1, 2) \\
$\alpha$ & User-specified risk level \\
$\delta$ & Significance level (default $0.05$) \\
$K$ & Sampling size for black-box UQ \\
$\mathrm{Coverage}(t)$ & Coverage: $|\{i : U_i \leq t\}| / N_{\mathrm{test}}$ \\
\bottomrule
\end{tabular}
\end{table}

\subsection{Implementation}
\label{app:implementation}

\paragraph{Candidate generation.}
We generate candidate answers using greedy decoding for both Qwen3-8B and LLaMA-3.1-70B. For Qwen3-8B, thinking mode is disabled (\texttt{enable\_thinking=False}) to prevent the \texttt{<think>} block from consuming the generation budget.

\paragraph{Judge inference.}
For predictive entropy (PE), we perform a single forward pass at the verdict token position and extract the log-probabilities for the \texttt{True} and \texttt{False} tokens. PE is computed as binary entropy: $-p_{\text{true}} \log p_{\text{true}} - (1 - p_{\text{true}}) \log (1 - p_{\text{true}})$. This requires no additional inference beyond the greedy decoding pass used to produce the verdict.

\paragraph{Retrieval.}
We query the Serper API with the original question and retrieve the top-$k$ results. Each result contains a title, text snippet, and source URL. Results are concatenated into a single evidence passage appended to the Mode~2 judge prompt.

\paragraph{Calibration.}
We run 100 random 50/50 calibration/test splits. For each split, we perform a grid search over all unique uncertainty values on the calibration set to find the threshold pair $(\hat{t}_1, \hat{t}_2)$ that maximizes correctly accepted instances subject to the Clopper--Pearson UCB constraint $\hat{R}^{\mathrm{upper}}_{1-\delta}(t_1, t_2) \leq \alpha$. The grid search uses cumulative sums for $O(|\mathcal{T}_1| \cdot |\mathcal{T}_2|)$ complexity.

\subsection{Prompt Templates}
\label{app:prompts}

\subsubsection{Judge Prompts}
\paragraph{Mode~1 (parametric evaluation).}
The judge receives a system prompt and a user message containing the question and candidate answer. The prompt includes a one-shot example to establish the output format.

\begin{tcolorbox}[colback=gray!5, colframe=gray!60, title={\small Mode 1: Judge System Prompt}, fonttitle=\bfseries\small, fontupper=\small]
You are a helpful assistant acting as an impartial judge. You will be given a Question and a Candidate Answer. Your task is to judge whether the Candidate Answer is correct using your own knowledge. If the Candidate Answer is correct, choose `True', otherwise, choose `False'. Provide a brief explanation for your decision.

\medskip
Provide your response in the following format:\\
\texttt{Decision: [True/False]}\\
\texttt{Explanation: [Your brief explanation]}

\medskip
\textit{Example:}\\
Question: What is the capital of France?\\
Candidate Answer: Paris is the capital of France.\\
Decision: True\\
Explanation: Paris is indeed the capital of France.
\end{tcolorbox}

\begin{tcolorbox}[colback=blue!3, colframe=blue!40, title={\small Mode 1: User Message}, fonttitle=\bfseries\small, fontupper=\small]
Question: \texttt{\{question\}}

Candidate Answer: \texttt{\{candidate\_answer\}}
\end{tcolorbox}

\paragraph{Mode~2 (retrieval-augmented evaluation).}
The judge additionally receives retrieved web search results. Each result is formatted as: \texttt{[i] Title / Snippet / Source: URL}.

\begin{tcolorbox}[colback=gray!5, colframe=gray!60, title={\small Mode 2: Judge System Prompt}, fonttitle=\bfseries\small, fontupper=\small]
You are a helpful assistant acting as an impartial judge. You will be given a Question, a Candidate Answer, and Web Search Results. Your task is to judge whether the Candidate Answer is correct using the provided search results. If the search results support the Candidate Answer, choose `True'. If the search results contradict it, choose `False'. If the search results are inconclusive, use your best judgment. Provide a brief explanation for your decision.

\medskip
Provide your response in the following format:\\
\texttt{Decision: [True/False]}\\
\texttt{Explanation: [Your brief explanation]}
\end{tcolorbox}

\begin{tcolorbox}[colback=blue!3, colframe=blue!40, title={\small Mode 2: User Message}, fonttitle=\bfseries\small, fontupper=\small]
Question: \texttt{\{question\}}

Candidate Answer: \texttt{\{candidate\_answer\}}

Web Search Results: \texttt{\{evidence\}}
\end{tcolorbox}


\subsection{LLM Evaluator Prompt}

The LLM evaluator (Qwen2.5-7B-Instruct) determines whether the candidate answer is semantically equivalent to any reference answer~\citep{badshah2025ref}, serving as the admission function for computing the ground-truth label~$e^*$.

\begin{tcolorbox}[colback=green!3, colframe=green!50!black, title={\small LLM Evaluator (Admission Function)}, fonttitle=\bfseries\small, fontupper=\small]
You are a helpful assistant acting as an impartial evaluator. Given a question, a candidate answer, and reference answers, determine if the candidate answer is correct.

\medskip
The candidate answer is correct if it conveys the same meaning as ANY of the reference answers, even if the candidate answer is longer, contains additional explanation, uses synonyms, abbreviations, or alternate names. Focus on whether the final answer or conclusion matches a reference, not on surface form.

\medskip
Question: \texttt{\{question\}}

Candidate Answer: \texttt{\{candidate\_answer\}}

Reference Answers: \texttt{\{references\}}

\medskip
Is the candidate answer correct? Respond with ONLY ``True'' or ``False''.
\end{tcolorbox}

\section{Additional Results}
\label{app:additional_results}
This appendix provides additional results that complement the main paper, including raw judge accuracy, an analysis of the joint framework across different risk levels, and per-dataset coverage and FDR curves.

\subsection{Judge Verdict Accuracy}
\label{app:accuracy}
Table~\ref{tab:verdict_accuracy} reports the raw verdict accuracy of each judge in Mode~1 (parametric) and Mode~2 (retrieval-augmented) before any calibration or selective prediction. Retrieval consistently improves accuracy across all configurations, with the largest gains on PopQA (up to 44 percentage points), showing that retrieval provides useful evidence.

\begin{table*}[!ht]
\centering
\small
\caption{Judge verdict accuracy (\%) before calibration. Mode~1 uses parametric knowledge only whereas Mode~2 uses retrieval-augmented evaluation ($k\!=\!3$). LLM evaluator as admission function.}
\label{tab:verdict_accuracy}
\setlength{\tabcolsep}{11.5pt}
\begin{tabular}{ll|cc|cc|cc|cc}
\toprule
 & & \multicolumn{2}{c|}{TriviaQA} & \multicolumn{2}{c|}{NQ-Open} & \multicolumn{2}{c|}{PopQA} & \multicolumn{2}{c}{HotpotQA} \\
Cand. & Judge & M1 & M2 & M1 & M2 & M1 & M2 & M1 & M2 \\
\midrule
\multirow{4}{*}{\rotatebox{90}{\scriptsize Qwen-8B}}
 & Qwen-4B  & 68 & 79 & 52 & 71 & 47 & 84 & 54 & 69 \\
 & Qwen-8B  & 71 & 84 & 54 & 79 & 44 & 88 & 58 & 76 \\
 & Qwen-14B & 78 & 86 & 61 & 77 & 65 & 88 & 67 & 78 \\
 & Llama-8B & 77 & 82 & 63 & 78 & 68 & 88 & 68 & 78 \\
\midrule
\multirow{4}{*}{\rotatebox{90}{\scriptsize Llama-70B}}
 & Qwen-4B  & 75 & 81 & 62 & 69 & 60 & 75 & 63 & 68 \\
 & Qwen-8B  & 78 & 81 & 63 & 71 & 61 & 77 & 65 & 69 \\
 & Qwen-14B & 78 & 82 & 64 & 72 & 66 & 77 & 67 & 72 \\
 & Llama-8B & 72 & 78 & 63 & 70 & 64 & 77 & 64 & 68 \\
\bottomrule
\end{tabular}
\end{table*}

\subsection{FDR and Coverage Across Risk Levels}
\label{app:full_results}
Table~\ref{tab:full_results} reports FDR and coverage for the joint two-mode framework (two-threshold routing with retrieval) at all evaluated $\alpha$ levels for the Qwen-8B candidate across all four datasets and all four judges (LLM evaluator, PE, $k\!=\!3$).

\begin{table*}[h]
\centering
\small
\caption{Full results for the joint two-mode framework at all $\alpha$ levels across all four datasets. Candidate: Qwen3-8B, LLM evaluator, PE, $k\!=\!3$.}
\label{tab:full_results}
\resizebox{\textwidth}{!}{%
\begin{tabular}{ll|cccccccccccc}
\toprule
 & & \multicolumn{12}{c}{$\alpha$} \\
Judge & Metric & .05 & .10 & .15 & .20 & .25 & .30 & .35 & .40 & .45 & .50 & .55 & .60 \\
\midrule
\multicolumn{14}{c}{\textit{TriviaQA}} \\
\midrule
\multirow{2}{*}{Qwen-4B}
 & FDR & .03 & .08 & .13 & .18 & .21 & .21 & .21 & .21 & .21 & .21 & .21 & .21 \\
 & Cov. & .03 & .39 & .66 & .91 & 1.0 & 1.0 & 1.0 & 1.0 & 1.0 & 1.0 & 1.0 & 1.0 \\
\multirow{2}{*}{Qwen-8B}
 & FDR & .03 & .08 & .14 & .16 & .16 & .16 & .16 & .16 & .16 & .16 & .16 & .16 \\
 & Cov. & .12 & .63 & .91 & 1.0 & 1.0 & 1.0 & 1.0 & 1.0 & 1.0 & 1.0 & 1.0 & 1.0 \\
\multirow{2}{*}{Qwen-14B}
 & FDR & .04 & .09 & .13 & .14 & .14 & .14 & .14 & .14 & .14 & .14 & .14 & .14 \\
 & Cov. & .21 & .78 & .98 & 1.0 & 1.0 & 1.0 & 1.0 & 1.0 & 1.0 & 1.0 & 1.0 & 1.0 \\
\multirow{2}{*}{Llama-8B}
 & FDR & .03 & .09 & .14 & .18 & .18 & .18 & .18 & .18 & .18 & .18 & .18 & .18 \\
 & Cov. & .13 & .61 & .86 & 1.0 & 1.0 & 1.0 & 1.0 & 1.0 & 1.0 & 1.0 & 1.0 & 1.0 \\
\midrule
\multicolumn{14}{c}{\textit{NQ-Open}} \\
\midrule
\multirow{2}{*}{Qwen-4B}
 & FDR & .00 & .05 & .11 & .17 & .22 & .27 & .29 & .29 & .29 & .29 & .29 & .29 \\
 & Cov. & .00 & .03 & .08 & .14 & .28 & .93 & 1.0 & 1.0 & 1.0 & 1.0 & 1.0 & 1.0 \\
\multirow{2}{*}{Qwen-8B}
 & FDR & .00 & .05 & .13 & .18 & .22 & .22 & .22 & .22 & .22 & .22 & .22 & .22 \\
 & Cov. & .00 & .02 & .24 & .82 & 1.0 & 1.0 & 1.0 & 1.0 & 1.0 & 1.0 & 1.0 & 1.0 \\
\multirow{2}{*}{Qwen-14B}
 & FDR & .00 & .05 & .13 & .18 & .22 & .23 & .23 & .23 & .23 & .23 & .23 & .23 \\
 & Cov. & .00 & .05 & .56 & .81 & .99 & 1.0 & 1.0 & 1.0 & 1.0 & 1.0 & 1.0 & 1.0 \\
\multirow{2}{*}{Llama-8B}
 & FDR & .00 & .09 & .13 & .18 & .22 & .22 & .22 & .22 & .22 & .22 & .22 & .22 \\
 & Cov. & .00 & .20 & .61 & .83 & 1.0 & 1.0 & 1.0 & 1.0 & 1.0 & 1.0 & 1.0 & 1.0 \\
\midrule
\multicolumn{14}{c}{\textit{PopQA}} \\
\midrule
\multirow{2}{*}{Qwen-4B}
 & FDR & .00 & .01 & .13 & .16 & .16 & .16 & .16 & .16 & .16 & .16 & .16 & .16 \\
 & Cov. & .00 & .01 & .91 & 1.0 & 1.0 & 1.0 & 1.0 & 1.0 & 1.0 & 1.0 & 1.0 & 1.0 \\
\multirow{2}{*}{Qwen-8B}
 & FDR & .00 & .09 & .12 & .12 & .12 & .12 & .12 & .12 & .12 & .12 & .12 & .12 \\
 & Cov. & .00 & .87 & 1.0 & 1.0 & 1.0 & 1.0 & 1.0 & 1.0 & 1.0 & 1.0 & 1.0 & 1.0 \\
\multirow{2}{*}{Qwen-14B}
 & FDR & .02 & .09 & .12 & .12 & .12 & .12 & .12 & .12 & .12 & .12 & .12 & .12 \\
 & Cov. & .19 & .89 & 1.0 & 1.0 & 1.0 & 1.0 & 1.0 & 1.0 & 1.0 & 1.0 & 1.0 & 1.0 \\
\multirow{2}{*}{Llama-8B}
 & FDR & .00 & .09 & .12 & .13 & .13 & .13 & .13 & .13 & .13 & .13 & .13 & .13 \\
 & Cov. & .02 & .88 & 1.0 & 1.0 & 1.0 & 1.0 & 1.0 & 1.0 & 1.0 & 1.0 & 1.0 & 1.0 \\
\midrule
\multicolumn{14}{c}{\textit{HotpotQA}} \\
\midrule
\multirow{2}{*}{Qwen-4B}
 & FDR & .00 & .05 & .10 & .15 & .21 & .27 & .31 & .31 & .31 & .31 & .31 & .31 \\
 & Cov. & .00 & .04 & .09 & .14 & .30 & .58 & .86 & .86 & .86 & .86 & .86 & .86 \\
\multirow{2}{*}{Qwen-8B}
 & FDR & .00 & .05 & .13 & .18 & .23 & .24 & .24 & .24 & .24 & .24 & .24 & .24 \\
 & Cov. & .00 & .04 & .49 & .74 & .96 & 1.0 & 1.0 & 1.0 & 1.0 & 1.0 & 1.0 & 1.0 \\
\multirow{2}{*}{Qwen-14B}
 & FDR & .02 & .08 & .13 & .18 & .22 & .22 & .22 & .22 & .22 & .22 & .22 & .22 \\
 & Cov. & .02 & .26 & .64 & .86 & 1.0 & 1.0 & 1.0 & 1.0 & 1.0 & 1.0 & 1.0 & 1.0 \\
\multirow{2}{*}{Llama-8B}
 & FDR & .01 & .08 & .13 & .18 & .22 & .22 & .22 & .22 & .22 & .22 & .22 & .22 \\
 & Cov. & .01 & .30 & .65 & .85 & 1.0 & 1.0 & 1.0 & 1.0 & 1.0 & 1.0 & 1.0 & 1.0 \\
\bottomrule
\end{tabular}%
}
\end{table*}

% \subsection{Additional Figures}
% \label{app:figures}
\subsubsection{Coverage Curves}
Figure~\ref{fig:coverage_all} shows coverage as a function of~$\alpha$ across all 32 configurations.

\begin{figure*}[!ht]
\centering
\includegraphics[width=\textwidth]{figures/coverage_all.pdf}
\caption{Coverage (mean$\pm$std over 100 splits) across all 32 configurations. Rows: candidates; columns: datasets.}
\label{fig:coverage_all}
\end{figure*}

\subsubsection{Per-Dataset FDR}
Figures~\ref{fig:fdr_triviaqa}--\ref{fig:fdr_hotpotqa} provide per-dataset FDR plots for detailed examination.

\begin{figure*}[!ht]
\centering
\includegraphics[width=0.85\textwidth]{figures/fdr_triviaqa.pdf}
\caption{FDR control on TriviaQA.}
\label{fig:fdr_triviaqa}
\end{figure*}

\begin{figure*}[!ht]
\centering
\includegraphics[width=0.85\textwidth]{figures/fdr_nqopen.pdf}
\caption{FDR control on NQ-Open.}
\label{fig:fdr_nqopen}
\end{figure*}

\begin{figure*}[!ht]
\centering
\includegraphics[width=0.85\textwidth]{figures/fdr_popqa.pdf}
\caption{FDR control on PopQA.}
\label{fig:fdr_popqa}
\end{figure*}

\begin{figure*}[!ht]
\centering
\includegraphics[width=0.85\textwidth]{figures/fdr_hotpotqa.pdf}
\caption{FDR control on HotpotQA.}
\label{fig:fdr_hotpotqa}
\end{figure*}

\section{Ablations and Robustness}
\label{app:ablations}
We evaluate the framework across several dimensions. This includes pipeline design choices, the calibration procedure, the uncertainty score, and the choice of no-retrieval baseline. In all cases, the finite-sample FDR guarantee is preserved, with conservative choices reducing coverage rather than compromising risk control.

\subsection{Sensitivity Analyses}
\label{sec:sensitivity}
\label{app:calibsens}

\paragraph{Correctness criterion (admission function).}
Table~\ref{tab:admission_sensitivity} compares three admission functions across three risk levels on a representative pair (Qwen3-8B candidate, Qwen3-14B judge). The LLM evaluator consistently achieves the highest coverage, with the gap most pronounced on knowledge-intensive benchmarks: on NQ-Open at $\alpha\!=\!0.15$, it reaches 56\% while both exact match and token~F1 remain at 0\%. This gap arises because string-matching criteria penalize semantically correct but surface-different answers, inflating the judge's apparent error rate. FDR control holds under all three criteria.

\begin{table*}[h]
\centering
\caption{Sensitivity to admission function. Candidate: Qwen3-8B, Judge: Qwen3-14B, UQ: PE (white-box). FDR and coverage at $\alpha \in \{.15, .20, .25\}$.}
\label{tab:admission_sensitivity}
\resizebox{\textwidth}{!}{%
\footnotesize
\begin{tabular}{l|ccc|ccc|ccc|ccc|ccc|ccc}
\toprule
 & \multicolumn{6}{c|}{Exact Match} & \multicolumn{6}{c|}{Token F1} & \multicolumn{6}{c}{LLM Evaluator} \\
 & \multicolumn{3}{c|}{FDR} & \multicolumn{3}{c|}{Coverage} & \multicolumn{3}{c|}{FDR} & \multicolumn{3}{c|}{Coverage} & \multicolumn{3}{c|}{FDR} & \multicolumn{3}{c}{Coverage} \\
Dataset & .15 & .20 & .25 & .15 & .20 & .25 & .15 & .20 & .25 & .15 & .20 & .25 & .15 & .20 & .25 & .15 & .20 & .25 \\
\midrule
TriviaQA & .13 & .14 & .14 & .97 & 1.0 & 1.0 & .13 & .18 & .18 & .72 & .98 & 1.0 & .13 & .14 & .14 & .98 & 1.0 & 1.0 \\
NQ-Open  & .00 & .05 & .22 & .00 & .11 & .69 & .00 & .00 & .18 & .00 & .00 & .45 & .13 & .18 & .22 & .56 & .81 & .99 \\
PopQA    & .12 & .12 & .12 & 1.0 & 1.0 & 1.0 & .13 & .13 & .13 & .99 & 1.0 & 1.0 & .12 & .12 & .12 & 1.0 & 1.0 & 1.0 \\
HotpotQA & .04 & .18 & .23 & .03 & .57 & .86 & .04 & .17 & .23 & .04 & .47 & .83 & .13 & .18 & .22 & .64 & .86 & 1.0 \\
\bottomrule
\end{tabular}%
}
\end{table*}

\paragraph{Retrieval depth.}
Table~\ref{tab:retrieval_depth} ablates the number of retrieved passages ($k \in \{3, 6, 9\}$) on a representative pair (Llama-70B candidate, Qwen3-8B judge). Coverage and FDR are stable across retrieval depths, with differences within 1--2 percentage points. $k\!=\!3$ already provides sufficient evidence.

\begin{table}[h]
\centering
\small
\caption{Retrieval depth ablation at $\alpha\!=\!0.20$. Candidate: Llama-70B, Judge: Qwen3-8B, LLM evaluator.}
\label{tab:retrieval_depth}
\setlength{\tabcolsep}{11.5pt}
\begin{tabular}{l|cc|cc|cc}
\toprule
 & \multicolumn{2}{c|}{$k=3$} & \multicolumn{2}{c|}{$k=6$} & \multicolumn{2}{c}{$k=9$} \\
Dataset & FDR & Cov. & FDR & Cov. & FDR & Cov. \\
\midrule
TriviaQA & .18 & .99 & .17 & 1.0 & .17 & .99 \\
NQ-Open  & .18 & .47 & .18 & .46 & .18 & .46 \\
PopQA    & .18 & .48 & .18 & .41 & .17 & .36 \\
HotpotQA & .18 & .47 & .18 & .46 & .18 & .46 \\
\bottomrule
\end{tabular}
\end{table}

\paragraph{Confidence level $\delta$.} Table~\ref{tab:delta} varies the confidence level. Tightening $\delta$ from $0.05$ to $0.01$ trades coverage for confidence, with the cost concentrated at low~$\alpha$ and on the harder dataset; FDR stays controlled throughout.

\begin{table}[h]
\centering
\small
\caption{Sensitivity to the confidence level $\delta$ (coverage). Candidate: Qwen3-8B, Judge: Qwen3-8B, LLM evaluator, $k\!=\!3$, PE. FDR is controlled at the target~$\alpha$ under both settings.}
\label{tab:delta}
\begin{tabular*}{\textwidth}{@{\extracolsep{\fill}}c|ccccc|ccccc@{}}
\toprule
 & \multicolumn{5}{c|}{TriviaQA} & \multicolumn{5}{c}{HotpotQA} \\
$\delta~\backslash~\alpha$ & .05 & .10 & .15 & .20 & .25 & .05 & .10 & .15 & .20 & .25 \\
\midrule
0.05 & .12 & .63 & .91 & 1.0 & 1.0 & .00 & .04 & .49 & .74 & .96 \\
0.01 & .04 & .53 & .87 & 1.0 & 1.0 & .00 & .01 & .37 & .69 & .93 \\
\bottomrule
\end{tabular*}
\end{table}

\paragraph{Calibration-set size.} Table~\ref{tab:calibsize} sweeps the calibration fraction $n_{\mathrm{cal}}/n$ over the 2{,}000 items per dataset. Coverage on the easier TriviaQA is stable across the range, while on HotpotQA it degrades gracefully as the calibration set shrinks; FDR is controlled at every size.

\begin{table}[h]
\centering
\small
\caption{Sensitivity to calibration-set size (coverage). $n_{\mathrm{cal}}/n$ is the fraction of the 2{,}000 items used for calibration. Candidate: Qwen3-8B, Judge: Qwen3-8B, LLM evaluator, $k\!=\!3$, PE. FDR is controlled at the target~$\alpha$ for every size.}
\label{tab:calibsize}
\begin{tabular*}{\textwidth}{@{\extracolsep{\fill}}c|ccccc|ccccc@{}}
\toprule
 & \multicolumn{5}{c|}{TriviaQA} & \multicolumn{5}{c}{HotpotQA} \\
$n_{\mathrm{cal}}/n~\backslash~\alpha$ & .05 & .10 & .15 & .20 & .25 & .05 & .10 & .15 & .20 & .25 \\
\midrule
0.10 & .02 & .38 & .81 & .96 & .99 & .00 & .03 & .25 & .55 & .83 \\
0.25 & .08 & .54 & .89 & .99 & 1.0 & .00 & .04 & .41 & .70 & .92 \\
0.50 & .12 & .63 & .91 & 1.0 & 1.0 & .00 & .04 & .49 & .74 & .96 \\
0.75 & .15 & .66 & .92 & 1.0 & 1.0 & .00 & .04 & .53 & .77 & .98 \\
\bottomrule
\end{tabular*}
\end{table}

\subsection{Alternative Uncertainty Estimators}
\label{app:uqmenu}
The FDR guarantee depends only on the binary correctness indicator. As a result, the choice of uncertainty score affects coverage but not whether the target risk level is controlled. In addition to the white-box predictive entropy (PE) used in the main paper, we evaluate two supervised uncertainty scores derived from the judge's last-layer hidden state at the verdict position. Following~\citet{NEURIPS2024_9c20f16b}, we consider a small MLP (\emph{Probe}) and a \emph{LoRA+Prompt} adapter fine-tuned to predict verdict correctness.

Both models are trained using 5-fold cross-validation on the calibration set so that every calibration instance receives an out-of-sample uncertainty score before conformal calibration. Table~\ref{tab:uqmenu} reports the resulting coverage. All three uncertainty scores satisfy the target FDR level across every setting.

LoRA+Prompt consistently achieves the highest coverage, particularly on datasets where PE is less informative. For example, coverage increases from $0.04$ to $0.44$ on HotpotQA at $\alpha\!=\!0.10$ and from $0.00$ to $0.71$ on PopQA at $\alpha\!=\!0.05$. The Probe also improves over PE on more challenging datasets at moderate risk levels (e.g., NQ-Open at $\alpha\!=\!0.15$, $0.24 \rightarrow 0.42$), although it performs worse at the lowest risk levels.

\begin{table}[!ht]
\centering
\small
\caption{Coverage under three uncertainty scores: the white-box predictive entropy (PE) used in the main paper and two supervised scores, a \emph{Probe} and a \emph{LoRA+Prompt} adapter on the judge's hidden states. Candidate: Qwen3-8B, Judge: Qwen3-8B, LLM evaluator, $k\!=\!3$. FDR is controlled at the target~$\alpha$ for all three.}
\label{tab:uqmenu}
\setlength{\tabcolsep}{4pt}
\begin{tabular*}{\textwidth}{@{\extracolsep{\fill}}c|ccc|ccc|ccc|ccc@{}}
\toprule
 & \multicolumn{3}{c|}{TriviaQA} & \multicolumn{3}{c|}{NQ-Open} & \multicolumn{3}{c|}{PopQA} & \multicolumn{3}{c}{HotpotQA} \\
$\alpha$ & PE & Probe & LoRA & PE & Probe & LoRA & PE & Probe & LoRA & PE & Probe & LoRA \\
\midrule
.05 & .12 & .00 & .00 & .00 & .00 & .02 & .00 & .31 & .71 & .00 & .00 & .02 \\
.10 & .63 & .14 & .70 & .02 & .07 & .28 & .87 & .82 & .94 & .04 & .02 & .44 \\
.15 & .91 & .90 & .94 & .24 & .42 & .64 & 1.0 & 1.0 & 1.0 & .49 & .47 & .76 \\
.20 & 1.0 & 1.0 & 1.0 & .82 & .79 & .86 & 1.0 & 1.0 & 1.0 & .74 & .78 & .95 \\
.25 & 1.0 & 1.0 & 1.0 & 1.0 & 1.0 & 1.0 & 1.0 & 1.0 & 1.0 & .96 & .99 & 1.0 \\
\bottomrule
\end{tabular*}
\end{table}

\subsection{Baseline Comparison}
\label{app:baselines}
To isolate the contribution of retrieved evidence, Table~\ref{tab:baselines} compares our joint two-mode policy against three controls that keep the identical calibration but change what the accepted verdict is based on. \emph{Direct} accepts only Mode~1 (no retrieval). \emph{Self-eval} adds a second Mode~1 pass in which the judge is asked whether its own initial verdict is correct, using no new evidence. \emph{Empty} runs the Mode~2 prompt with an empty evidence slot. FDR is controlled for every variant. The joint policy dominates all three at nearly every operating point, and the gap is largest on the knowledge-intensive HotpotQA ($0.74$ versus at most $0.07$ for the no-retrieval controls at $\alpha\!=\!0.20$). A second evaluation pass without new evidence (Self-eval) barely helps, confirming that the coverage gains come from the retrieved evidence rather than from simply re-querying the judge.

\begin{table}[!ht]
\centering
\small
\caption{Baseline comparison (coverage). Candidate: Qwen3-8B, Judge: Qwen3-8B, LLM evaluator, $k\!=\!3$, PE. FDR is controlled at the target~$\alpha$ for every variant.}
\label{tab:baselines}
\begin{tabular*}{\textwidth}{@{\extracolsep{\fill}}l|ccccc|ccccc@{}}
\toprule
 & \multicolumn{5}{c|}{TriviaQA} & \multicolumn{5}{c}{HotpotQA} \\
Variant~$\backslash$~$\alpha$ & .05 & .10 & .15 & .20 & .25 & .05 & .10 & .15 & .20 & .25 \\
\midrule
Direct       & .03 & .28 & .47 & .59 & .74 & .00 & .00 & .01 & .05 & .17 \\
Self-eval    & .02 & .32 & .49 & .60 & .75 & .00 & .00 & .02 & .07 & .19 \\
Empty        & .04 & .38 & .63 & .83 & .99 & .00 & .09 & .32 & .57 & .84 \\
\midrule
Joint (ours) & .12 & .63 & .91 & 1.0 & 1.0 & .00 & .04 & .49 & .74 & .96 \\
\bottomrule
\end{tabular*}
\end{table}

\subsection{Conservative Joint Calibration}
\label{app:bonferroni}
We selected the deployed threshold pair $(\hat{t}_1, \hat{t}_2)$ using the calibration data rather than fixing them beforehand. Specifically, we choose the pair that maximizes coverage while satisfying $\hat{R}^{\mathrm{upper}}_{1-\delta} \leq \alpha$ (Eq.~\ref{eq:that_route}). Although the Clopper--Pearson bound holds for each candidate pair on its own, certifying the selected pair requires accounting for this search. The conservative variant addresses this directly by evaluating all $|\mathcal{T}_1|\,|\mathcal{T}_2|$ candidate pairs using the corrected confidence level $\delta'=\delta/(|\mathcal{T}_1|\,|\mathcal{T}_2|)$. By a union-bound argument, the guarantee then holds simultaneously for every candidate pair, implying that the selected pair satisfies $\Pr(R(\hat{t}_1,\hat{t}_2)\leq\alpha)\geq1-\delta$. In practice, however, we find that the simpler uncorrected selection already provides effective risk control (see Section~\ref{sec:results}). For this reason, we use the coverage-maximizing thresholds in our main experiments, while treating the conservative procedure as a theoretically valid alternative.

One subtlety is that the grid is data-dependent, since $\mathcal{T}_1$ and $\mathcal{T}_2$ are the observed uncertainty values, making both the candidate pairs and their count random. Because each mode contributes at most $N$ distinct thresholds, replacing $|\mathcal{T}_1|\,|\mathcal{T}_2|$ with the data-independent bound $N^2$ (equivalently, correcting at level $\delta/N^2$) makes the union bound hold over a fixed family and leaves the conclusion unchanged at negligible cost; pre-specifying the grid removes the dependence altogether.

Table~\ref{tab:bonferroni} compares the coverage of the main-paper calibration (\emph{Pointwise}, per-pair level $\delta$) with this \emph{Conservative} variant (grid-corrected level $\delta'$). FDR is controlled under both. The correction is nearly free on the easier datasets (TriviaQA, PopQA) but costs substantial coverage on the harder ones (NQ-Open, HotpotQA), where the empirical FDR sits close to~$\alpha$. Bonferroni is a worst-case union bound and adjacent grid thresholds select almost the same instances, so it is very loose. The truly attainable jointly-valid coverage lies between the two rows which is much closer to the pointwise figures.

\begin{table}[h]
\centering
\small
\caption{Conservative joint calibration (coverage). Per dataset, \emph{Point.}\ and \emph{Cons.}\ denote the pointwise and conservative coverage, and $\Delta = \text{Cons.} - \text{Point.}$ (computed on unrounded values). Candidate: Qwen3-8B, Judge: Qwen3-8B, LLM evaluator, $k\!=\!3$, PE. FDR is controlled at the target~$\alpha$ under both variants.}
\label{tab:bonferroni}
\setlength{\tabcolsep}{4pt}
\begin{tabular*}{\textwidth}{@{\extracolsep{\fill}}c|ccc|ccc|ccc|ccc@{}}
\toprule
 & \multicolumn{3}{c|}{TriviaQA} & \multicolumn{3}{c|}{NQ-Open} & \multicolumn{3}{c|}{PopQA} & \multicolumn{3}{c}{HotpotQA} \\
$\alpha$ & Point. & Cons. & $\Delta$ & Point. & Cons. & $\Delta$ & Point. & Cons. & $\Delta$ & Point. & Cons. & $\Delta$ \\
\midrule
.05 & .12 & .00 & $-$.12 & .00 & .00 & .00 & .00 & .00 & .00 & .00 & .00 & .00 \\
.10 & .63 & .01 & $-$.62 & .02 & .00 & $-$.02 & .87 & .00 & $-$.87 & .04 & .00 & $-$.04 \\
.15 & .91 & .64 & $-$.27 & .24 & .00 & $-$.24 & 1.0 & .93 & $-$.07 & .49 & .00 & $-$.49 \\
.20 & 1.0 & .93 & $-$.07 & .82 & .05 & $-$.76 & 1.0 & 1.0 & .00 & .74 & .30 & $-$.45 \\
.25 & 1.0 & 1.0 & .00 & 1.0 & .80 & $-$.20 & 1.0 & 1.0 & .00 & .96 & .72 & $-$.24 \\
\bottomrule
\end{tabular*}
\end{table}

% ==== Analysis ====
\section{Analysis}
\label{app:analysis}
This section moves beyond aggregate metrics to analyze how retrieval improves evaluation, including its impact on judge uncertainty, examples where it corrects incorrect verdicts, and the residual failure modes.

\subsection{Effect of Retrieval on Uncertainty}
\label{app:deltau}
Retrieval does not lower the judge's uncertainty uniformly. Table~\ref{tab:deltau} reports, per dataset, the fraction of items on which Mode~2 uncertainty is higher than Mode~1 ($U_2 > U_1$) versus lower, with the mean signed change $\Delta U = U_2 - U_1$ within each group. On most items retrieval reduces uncertainty, but on 14--29\% it raises it, concentrated on the multi-hop HotpotQA and NQ-Open, where retrieved snippets are only partially relevant or sources disagree. These are exactly the items the second threshold~$\hat{t}_2$ routes to abstention when retrieval fails to resolve the judge's uncertainty.

\begin{table}[!ht]
\centering
\small
\caption{Effect of retrieval on judge uncertainty. Results are obtained with Qwen3-8B as both candidate and judge, the LLM evaluator, and $k\!=\!3$ retrieval using predictive entropy (PE). The ``raised'' and ``lowered'' columns report the fraction of instances where retrieval increases or decreases uncertainty, respectively. Mean~$\Delta U$ is averaged within each group.}
\label{tab:deltau}
\begin{tabular*}{\textwidth}{@{\extracolsep{\fill}}l|cccc@{}}
\toprule
 & TriviaQA & NQ-Open & PopQA & HotpotQA \\
\midrule
\% raised ($U_2\!>\!U_1$)  & 20.1 & 27.9 & 14.3 & 29.1 \\
\quad mean $\Delta U$      & $+0.10$ & $+0.11$ & $+0.10$ & $+0.12$ \\
\% lowered ($U_2\!<\!U_1$) & 68.0 & 68.8 & 84.6 & 69.7 \\
\quad mean $\Delta U$      & $-0.07$ & $-0.09$ & $-0.10$ & $-0.12$ \\
\bottomrule
\end{tabular*}
\end{table}

\subsection{Qualitative Examples}
\label{app:qualitative}

Tables~\ref{tab:qualitative} and~\ref{tab:evidence_examples} show representative cases where the judge issues an incorrect verdict in Mode~1 but corrects it after retrieval in Mode~2. In each case, Mode~1 confidently hallucinated while retrieved evidence grounded the judge in verifiable facts.

\begin{table*}[!ht]
\centering
\caption{Examples where retrieval corrects the judge's verdict. Candidate: Qwen3-8B, Judge: Qwen3-14B.}
\label{tab:qualitative}
\resizebox{\textwidth}{!}{%
\small
\begin{tabular}{p{3.5cm}p{2cm}p{1.5cm}p{5.5cm}p{5.5cm}}
\toprule
\textbf{Question} & \textbf{Candidate} & \textbf{Gold} & \textbf{Mode 1 (Wrong)} & \textbf{Mode 2 (Corrected)} \\
\midrule
Who voices Jack Kahuna Laguna in SpongeBob SquarePants? & Steve Carell & Johnny Depp & \textit{``Steve Carell provided the voice\ldots This is a well-documented fact.''} \newline \textbf{Verdict: True} & \textit{``The web search results indicate that Johnny Depp, not Steve Carell, provided the voice\ldots''} \newline \textbf{Verdict: False} \\
\midrule
Which North African country's flag is just the colour green? & Chad & Libya & \textit{``The flag of Chad is indeed a single solid green rectangle\ldots''} \newline \textbf{Verdict: True} & \textit{``The search results indicate that Libya's flag was a single green field\ldots''} \newline \textbf{Verdict: False} \\
\midrule
`The Great American Chocolate Bar' is better known as? & Milky Way & Hershey & \textit{``The Great American Chocolate Bar is indeed known as the Milky Way bar\ldots''} \newline \textbf{Verdict: True} & \textit{``The search results indicate that it is a nickname for the Hershey Milk Chocolate Bar\ldots''} \newline \textbf{Verdict: False} \\
\bottomrule
\end{tabular}%
}
\end{table*}

\begin{table*}[!ht]
\centering
\caption{Retrieved evidence ($k\!=\!3$) for the examples in Table~\ref{tab:qualitative}. The first result in each case directly contains the correct answer.}
\label{tab:evidence_examples}
\resizebox{\textwidth}{!}{%
\small
\begin{tabular}{p{3.5cm}cp{12cm}}
\toprule
\textbf{Question} & \textbf{\#} & \textbf{Retrieved Snippet} \\
\midrule
\multirow{3}{3.5cm}{Who voices Jack Kahuna Laguna in SpongeBob?}
 & 1 & \textit{``\ldots American actor and musician Johnny Depp guest starred in the episode as the voice of Jack Kahuna Laguna, a surf guru\ldots''} --- Wikipedia \\
 & 2 & \textit{``Jack Kahuna Laguna has a striking resemblance to his voice actor, Johnny Depp.''} --- Fandom \\
 & 3 & \textit{``\ldots it guest starred Johnny Depp as the legendary surfer, Jack Kahuna Laguna.''} --- IMDb \\
\midrule
\multirow{3}{3.5cm}{Which N.\ African country's flag is just green?}
 & 1 & \textit{``The flag of the Libyan Arab Jamahiriya\ldots consisted of a simple green field. It was the only national flag\ldots with only one colour.''} --- Wikipedia \\
 & 2 & \textit{``Republic of Ireland and the Ivory Coast the flags are the same\ldots''} --- Facebook \\
 & 3 & \textit{``Green is an important colour in Islam\ldots''} --- Reddit \\
\midrule
\multirow{3}{3.5cm}{`The Great American Chocolate Bar' is better known as?}
 & 1 & \textit{``The classic HERSHEY'S Milk Chocolate bar has been referred to as `The Great American Chocolate Bar.'\,''} --- Hershey \\
 & 2 & \textit{``Hershey refers to it as `The Great American Chocolate Bar'. First sold in 1900.''} --- Wikipedia \\
 & 3 & \textit{``The Great American Chocolate Bar campaign served the company well\ldots''} --- Hershey Archives \\
\bottomrule
\end{tabular}%
}
\end{table*}

\subsection{Failure Modes}
\label{app:failuremodes}
When retrieval does not fix an incorrect Mode~1 verdict, the residual errors fall into four qualitative classes. (i)~\emph{Off-topic retrieval}: an ambiguous or under-specified question yields search queries that miss the relevant evidence. (ii)~\emph{Insufficient evidence}: the retrieved snippets are only partially relevant and lack the specific fact needed to verify the candidate. (iii)~\emph{Judge reasoning error}: the snippets contain the relevant evidence, but the judge still returns a wrong verdict. For example defaulting to \texttt{False} in the absence of an exact lexical match. (iv)~\emph{Admission function disagreement}: the candidate is semantically close to the gold answer but the admission function scores it as different. The two-threshold policy sends many such cases to abstention rather than accepting them, which is why coverage rather than FDR absorbs the difficulty on the hardest datasets.

